\chapter{Experimental Setup}\label{ch:setup}

The evaluation has two phases. First, the generative models and projection methods are established on
the obstacle-avoidance benchmark of \ac{DPCC}, with its setup taken over unchanged so that the diffusion
baseline is met on its own benchmark. Second, they are tested further in two environments built for
this thesis: a vision-conditioned alignment task and a quadrotor benchmark with three scenes.
Throughout this chapter and the results, these environments are called \emph{D3IL-avoiding},
\emph{D3IL-aligning}, \emph{UAV-corridor}, \emph{UAV-pillars}, and \emph{UAV-s-curve}; the descriptive
subsection titles and these prefixed names refer to the same tasks. \autoref{sec:method:envs} describes how each environment is constructed; this
chapter gives, for each of the three, the task and constraint values, the data, the metrics it is scored
on, the baseline it is compared with, and the training and evaluation protocol.

\section{Simulated Platforms}\label{sec:setup:platforms}

Every environment is simulated by MuJoCo, the rigid-body physics engine used throughout this thesis
(\autoref{sec:bg:quadrotor}) \parencite{todorov2012mujoco}. The two
manipulation tasks run the seven-joint Panda manipulator in MuJoCo, with the model files, task
definitions and demonstrations of D3IL \parencite{jia2024towards}. The quadrotor benchmark runs the Skydio X2 model from MuJoCo Menagerie
\parencite{zakka2022menagerie} with the quadrotor task patch distributed by MuJoCo MPC
\parencite{howell2022predictive}. MuJoCo is the simulator in every case; D3IL and MuJoCo MPC supply the
task files and data that run in it. \autoref{fig:platforms} shows both models; the X2 dimensions used
in projection are drawn in \autoref{fig:platform-x2-dimensions} in the appendix.
\autoref{tab:platforms} gives the properties that enter the experiments.

How a plan reaches each of them differs. On the manipulator the planner commands an end-effector
increment; inverse kinematics turns each Cartesian setpoint into joint commands. On the quadrotor the
planner updates a commanded position increment once every three $0.01$\,s MuJoCo physics steps on
UAV-corridor and UAV-s-curve, and once every hundred on UAV-pillars, which flies the planar planner of
D3IL-avoiding at a fixed altitude (\autoref{sec:setup:protocol:uav}).
Between planner updates, the cascaded geometric controller of \autoref{sec:method:geometric}
computes rotor thrusts at every physics step.

\begin{figure}[htpb]
  \centering
  \begin{minipage}[b]{0.47\linewidth}\centering
    \includegraphics[height=2.8cm]{fig_platform_panda}\\[0.3em]
    {\footnotesize Panda manipulator}
  \end{minipage}\hfill
  \begin{minipage}[b]{0.47\linewidth}\centering
    \includegraphics[height=2.8cm]{fig_platform_x2}\\[0.3em]
    {\footnotesize Skydio X2 quadrotor}
  \end{minipage}
  \caption[The two simulated platforms]{The two simulated platforms, rendered in MuJoCo from the model
  files the experiments load, without their environments. Left: the Panda manipulator, from the model files
  of D3IL, used for D3IL-avoiding and D3IL-aligning. Right: the Skydio X2 model from MuJoCo
  Menagerie with the MuJoCo MPC task patch, used in the three aerial scenes. Both are simulated by
  MuJoCo. The two renders are not to scale with each other. The dimensions of both platforms are in
  \autoref{tab:platforms}; the quadrotor's rotor layout and the footprint the projection uses for it are
  drawn in \autoref{fig:platform-x2-dimensions} in the appendix.}
  \label{fig:platforms}
\end{figure}

\begin{table}[htpb]
  \caption[Properties of the simulated platforms]{Properties of the two platforms that enter the
  experiments. The Panda is taken from the model and task files of D3IL, the X2 from MuJoCo Menagerie
  with the MuJoCo MPC task patch; the X2 footprint includes the rotor disks of the collision model
  (\autoref{fig:platform-x2-dimensions} in the appendix).}
  \label{tab:platforms}
  \centering
  \footnotesize
  \setlength{\tabcolsep}{4pt}
  \begin{tabularx}{\linewidth}{@{}p{2.6cm} L p{2.5cm} p{2.5cm}@{}}
    \toprule
    Platform & Model & Plan commands & Plan interval \\
    \midrule
    Franka Emika Panda
      & seven joints; rod end effector, radius $0.01$\,m, length $0.30$\,m; wrist RGB camera
      & end-effector increments, planar or spatial
      & every physics step \\
    Skydio X2
      & four rotors; mass $1.325$\,kg; rotor centres $0.228$\,m out; $0.54\times0.62$\,m footprint;
        thrust $[0,13]$\,N per rotor
      & spatial position increments (UAV-corridor, UAV-s-curve); planar increments at a fixed altitude (UAV-pillars)
      & $0.03$\,s of simulated time, $1$\,s on UAV-pillars; controller at $0.01$\,s \\
    \bottomrule
  \end{tabularx}
\end{table}

\section{Tasks and Environments}\label{sec:setup:tasks}

The three environments differ in what the planner observes, how many dimensions its plan has, and which
constraint forms the projection must satisfy. Each is given below with its task, its constraint set and the
values that enter the experiments.

\subsection{Obstacle Avoidance}\label{sec:setup:tasks:avoiding}

The obstacle-avoidance task is the benchmark of \ac{DPCC}, simulated in MuJoCo from the task files and
demonstrations of D3IL: the same environment, Panda manipulator, demonstrations, constraint geometries, projection configuration and evaluation
protocol \parencite{jia2024towards,romer2025diffusion} (\autoref{sec:method:env:avoiding}). Holding
that complete benchmark fixed provides the controlled validation of replacing DPCC's diffusion model
with the three flow-based objectives. A manipulator moves its
end effector across a field of six obstacles to a goal line (\autoref{fig:env-avoiding}).

\begin{figure}[htpb]
  \centering
  \begin{minipage}[b]{0.47\linewidth}\centering
    \includegraphics[width=\linewidth]{fig_render_avoiding_start}\\[0.3em]
    {\footnotesize Start pose}
  \end{minipage}\hfill
  \begin{minipage}[b]{0.47\linewidth}\centering
    \includegraphics[width=\linewidth]{fig_render_avoiding}\\[0.3em]
    {\footnotesize Goal passed}
  \end{minipage}
  \caption[The obstacle-avoidance task in MuJoCo]{The Panda manipulator in the MuJoCo obstacle-avoidance
  scene. Left: the starting pose, with six red posts between the arm and the goal line. Right: the
  end effector beyond the green goal line.}
  \label{fig:env-avoiding}
\end{figure}

At test time
the plan is constrained by halfspaces that the demonstrations never respect, in three geometries --- \emph{top-left-hard}, \emph{top-right-hard} and \emph{both-hard}, defined in \autoref{tab:avoiding-geometries} and drawn in \autoref{fig:constraints-avoiding-sets} --- together with the circular keep-out regions
of \autoref{sec:method:env:avoiding}. The geometries are those of the \ac{DPCC} configuration \parencite[Sec.~6.1]{romer2025diffusion}, and all three share the workspace $x\in[0.2,0.8]$, $y\in[-0.3,0.4]$, the six obstacles --- one of radius $0.03$, five of $0.025$ --- and the goal line at $y=0.35$. Every
constraint also has a tightened version, shrunk by the margin of $0.025$ that \ac{DPCC} introduces and
applies to this benchmark \parencite[Sec.~5.3]{romer2025diffusion} (\autoref{sec:method:dpcc}). The same
holds for every constraint set of this thesis, and one convention is fixed here for every result that
follows: \textbf{every result of \autoref{ch:results} is on the tightened set.} It is the set \ac{DPCC}
reports on, and it is the set that gives the better result: on this benchmark, tightening raises the
published rate of success with constraint satisfaction from $0.49$ to $0.98$ and cuts the violating steps
from $6.0$ to $0.0$ per episode \parencite[Table~1]{romer2025diffusion}. \autoref{sec:setup:baselines:avoiding}
quotes those numbers in full next to their reproduction; the shrinkage itself, the Minkowski difference
with the mismatch ball, is stated in \autoref{sec:method:dpcc}. Untightened evaluations are therefore not
reported. The tightened set is the set the projector is handed; every violation count of
\autoref{ch:results} is taken against the nominal constraints, the boundary the tightening keeps a margin
from, and on UAV-corridor and UAV-s-curve the nominal surfaces carry the vehicle's rotor reach
(\autoref{sec:setup:metrics:uav}). An unprojected plan is the same plan on either set and is scored the
same way.

\begin{table}[htpb]
  \caption[Constraint geometries of obstacle avoidance]{The three constraint geometries of the
  obstacle-avoidance task, in metres. A halfspace keeps the planned end-effector position below the line
  through its two points and the keep-out disk outside the circle; the tightened version of each moves
  the line and enlarges the disk by $0.025$. The last column counts the D3IL demonstrations that satisfy
  the geometry at every step.}
  \label{tab:avoiding-geometries}
  \centering
  \begin{tabularx}{\linewidth}{@{}l L l r@{}}
    \toprule
    Geometry & Halfspace: below the line through & Keep-out disk & Satisfied by \\
    \midrule
    top-left-hard  & $(0.8,\,-0.5)$ and $(0.4,\,0.5)$ & $(0.4,\,0.08)$, $r=0.08$ & 0 of 96 \\
    top-right-hard & $(0.2,\,-0.5)$ and $(0.6,\,0.5)$ & $(0.6,\,0.08)$, $r=0.08$ & 1 of 96 \\
    both-hard      & $(0.8,\,-0.3)$ and $(0.575,\,0.5)$; $(0.2,\,-0.3)$ and $(0.425,\,0.5)$ & $(0.5,\,-0.09)$, $r=0.08$ & 2 of 96 \\
    \bottomrule
  \end{tabularx}
\end{table}

Each geometry is two things (\autoref{fig:constraints-avoiding-sets}). The \emph{halfspace} is a straight
line across the field with one side declared out of bounds for the planned end-effector position: the plan
must stay below the line, so the region above it --- the shaded corner --- is excluded, and the constraint
is a single linear row per plan step. The \emph{keep-out disk} is a circle of radius $0.08$\,m that the
planned position must stay outside of, placed over one of the posts so that it covers the post and a margin
of free table around it; it is the quadratic row of the set, and it is what makes the free lane through the
field a corridor between the halfspace on one side and the disk on the other. \emph{both-hard} has two
halfspaces, one from each side, and its disk on the first post. Neither element is a physical object: the
posts are the only obstacles the arm can collide with, and a plan that crosses a halfspace or enters a disk
violates the constraint without any collision. The tightened version of each set moves the line and enlarges the disk by
$0.025$\,m, the dashed boundaries of the figure.

\begin{figure}[htpb]
  \centering
  \includegraphics[width=\linewidth]{fig_constraints_avoiding_sets}
  \caption[The three constraint geometries of obstacle avoidance]{The three constraint geometries of
  \autoref{tab:avoiding-geometries}, seen from above on the field of \autoref{fig:env-avoiding}, with
  nothing but the constraints drawn: the shaded region is the side of the halfspace the planned position is
  excluded from, the shaded circle is the keep-out disk it must stay outside of, the dashed boundaries are the
  tightened versions, the red disks are the six posts and the green line the goal line. The demonstrations
  are drawn over the same three panels in \autoref{fig:constraints-avoiding}.}
  \label{fig:constraints-avoiding-sets}
\end{figure}

\FloatBarrier

\subsection{Vision-Conditioned Alignment}\label{sec:setup:tasks:aligning}

The alignment task of D3IL \parencite{jia2024towards} asks the manipulator to push a box into a target
pose on a table. The task runs in MuJoCo; its task files, demonstrations and camera set-up are D3IL's. Around them this thesis
built a constrained-control environment (\autoref{sec:method:env:aligning}): a three-dimensional plan
with its dynamics rows, a constraint set on the planned end-effector position, and the coupling of the
visual encoder to the generative network. The policy observes two RGB camera
streams, a front view and an in-hand view, instead of the state (\autoref{fig:aligning-cameras}). Each
episode starts from an initial box pose and a target pose, together called a context. The tightened
constraints are shrunk by $0.03$\,m.

Where the box is can only be seen. The vector half of the observation is
$(\vect{p}\sidx{des},\vect{p})\in\R^{6}$, the commanded and the measured end-effector position
(\autoref{sec:method:env:aligning}); it carries no box pose, no target pose and no object state of any
kind. The two RGB streams are the sole source of both, so the policy must locate the box in the images
before it can move toward it --- which is why a $0.10$\,m box in a $0.20\times0.15$\,m draw region is a
harder problem than its planar geometry suggests.

\begin{figure}[htpb]
  \centering
  \begin{minipage}{0.44\linewidth}\centering
    \includegraphics[width=\linewidth]{fig_render_aligning}\\[0.3em]
    {\footnotesize MuJoCo simulator view}
  \end{minipage}\hfill
  \begin{minipage}{0.44\linewidth}\centering
    \includegraphics[width=\linewidth]{fig_scene_aligning}\\[0.3em]
    {\footnotesize orthographic task-geometry view}
  \end{minipage}
  \caption[The vision-conditioned alignment task]{The vision-conditioned alignment task. Left: a frame
  rendered by MuJoCo with the manipulator, movable box and coloured target.
  Right: a geometry view of the first held-out context, reconstructed from the scene description
  and its recorded box and target poses. The transparent footprint marks the target pose. These panels
  show the task; the policy itself observes the front and in-hand RGB streams.}
  \label{fig:env-aligning}
\end{figure}

\begin{figure}[htpb]
  \centering
  \begin{minipage}{0.33\linewidth}\centering
    \includegraphics[width=\linewidth]{fig_aligning_camera_overhead}\\[0.3em]
    {\footnotesize front view}
  \end{minipage}\hspace{0.08\linewidth}
  \begin{minipage}{0.33\linewidth}\centering
    \includegraphics[width=\linewidth]{fig_aligning_camera_wrist}\\[0.3em]
    {\footnotesize in-hand view}
  \end{minipage}
  \caption[The two camera observations used for alignment]{The two $96\times96$ RGB observations the
  visual policy consumes, at one instant of a recorded demonstration with the box already at its target,
  rendered by MuJoCo. Left: the front view of manipulator, box and table. Right: the in-hand view of the
  box and its coloured sides. Camera placements are those of the D3IL task files
  \parencite{jia2024towards}.}
  \label{fig:aligning-cameras}
\end{figure}

Obstacle avoidance is a planar task that imitation policies already solve: the best policies reported
by D3IL reach a success rate of $0.95$ on it \parencite[Table~3]{jia2024towards}, so it tests constraint
satisfaction on a plan that is easy to produce. This environment raises three things at once. The plan is
three-dimensional rather than planar; the policy sees the scene through two cameras and a visual encoder
rather than through a state vector; and the task itself is not solved --- the best image-based policy
reported by D3IL finishes roughly a third of its episodes (\autoref{sec:setup:baselines:aligning},
\autoref{tab:va-vs-d3il}). It therefore asks whether the \ac{DPCC} construction still holds where the plan
is harder to produce, and it puts the projection where only some episodes finish at all.

The active geometry is \emph{combined-5} (\autoref{tab:aligning-constraints},
\autoref{fig:aligning-contexts}): a halfspace that cuts off the far corner of the table, a circular
keep-out region between the box and its target --- the same two geometric forms D3IL-avoiding uses, on a
task whose plan is three-dimensional. Neither is built into the demonstrations: the straight push from
the box to its target crosses the keep-out region in most recorded contexts and never the halfspace, and
the plans learned from the demonstrations enter both (\autoref{sec:setup:data:aligning},
\autoref{sec:res:aligning:projection}); satisfying the two is left to the projection.

\begin{table}[htpb]
  \caption[Constraint set of the alignment task]{The constraint set of the alignment task, geometry
  \emph{combined-5}, in metres, on the planned end-effector position. The tightened version moves the
  halfspace $0.03$\,m inward and grows the keep-out radius by the same margin.}
  \label{tab:aligning-constraints}
  \centering
  \small
  \setlength{\tabcolsep}{8pt}
  \renewcommand{\arraystretch}{1.2}
  \begin{tabularx}{\linewidth}{@{}l>{\raggedright\arraybackslash}p{3.9cm} X@{}}
    \toprule
    Constraint & Form & Values \\
    \midrule
    Halfspace & below the line through two points & $(0.65,\,0.45)$ and $(0.80,\,-0.45)$ \\
    Keep-out region & outside a circle & centre $(0.50,\,0.00)$, radius $0.06$ \\
    Tightening & both, moved inward & $0.03$ \\
    \bottomrule
  \end{tabularx}
\end{table}

The context is drawn at random, and this is the difference from obstacle avoidance that matters most for
what the policy has to do. In obstacle avoidance every episode begins at one fixed end-effector position
and the six posts never move, so a policy can learn a route. Here the box is placed anywhere in
$x\in[0.4,0.6]$, $y\in[-0.25,-0.10]$ at any yaw in $\pm 90^{\circ}$, and its target anywhere in
$x\in[0.4,0.6]$, $y\in[0.20,0.35]$ at any yaw in $\pm 90^{\circ}$; the demonstrations were recorded from
such draws and the evaluation replays an enumerated set of them (\autoref{sec:setup:protocol:aligning}).
Two episodes therefore ask for two different pushes. \autoref{fig:aligning-contexts} draws the ten
contexts the evaluation replays, which are the ten every result in \autoref{sec:res:visual} is averaged
over. They are ten different pushes: their lengths range from $0.41$ to $0.52$\,m and average
$0.4530$\,m, the starting distance \autoref{sec:res:visual} reports, so a policy that had learned one
route would satisfy none of them. One context in that figure is drawn with the real box footprints at
their recorded yaw, which every draw carries, because the marks alone understate the object --- a
$0.10$\,m box has to be pushed around a $0.06$\,m keep-out region, with the end effector kept out of the
far corner the halfspace excludes.

\begin{figure}[htpb]
  \centering
  \includegraphics[width=0.62\linewidth]{fig_aligning_contexts}
  \caption[The ten evaluated alignment contexts]{The ten contexts the evaluation replays, on the plane
  of \autoref{tab:aligning-constraints}. Filled marks are the box at the start of an episode, hollow
  marks its target, and a dotted line joins the two poses of one context; the dashed rectangles are the
  regions the environment draws each pose from. Both marks are box \emph{centres}, which is what the
  context and the final-distance metric refer to. One context also carries the real $0.10$\,m footprints
  at their recorded yaw.}
  \label{fig:aligning-contexts}
\end{figure}

\FloatBarrier

\subsection{Quadrotor Benchmark}\label{sec:setup:tasks:uav}

The quadrotor benchmark was built for this thesis around the Skydio X2 model of MuJoCo Menagerie
\parencite{zakka2022menagerie} (\autoref{sec:method:env:uav}). The vehicle is underactuated and open-loop
unstable (\autoref{sec:bg:quadrotor}), and its plan is tracked by the controller of
\autoref{sec:method:deployment}. Three scenes are defined
(\autoref{fig:env-uav}, \autoref{tab:uav-scenes}):

\begin{description}
  \item[UAV-pillars] --- D3IL-avoiding flown by the quadrotor. The obstacle field of
    \autoref{sec:setup:tasks:avoiding} is scaled by $36$, so that the rod's radius becomes the vehicle's
    rotor reach, and its six obstacles are extruded into pillars; the three test-time constraint
    geometries of that benchmark are mapped into the arena at the same scale. The planner is the one
    trained on the avoiding demonstrations and is not retrained: the quadrotor flies the plans it
    produces. The scene takes the configurations that already solve D3IL-avoiding, flies them,
    and reads the flown outcome beside the table outcome of the same configuration, before and after
    projection (\autoref{sec:res:uav:pillars}). The constraint set of each geometry is the table's: one or
    two halfspaces and one keep-out disk, which covers one of the six pillars. The other five pillars are
    obstacles of the scene and belong to no constraint set, on D3IL-avoiding or on UAV-pillars; a flight in
    which the vehicle collides with one ends there.
  \item[UAV-corridor] --- the three-dimensional test. Two parallel walls leave a clear width of
    $1.90$\,m, and the scene carries two test-time constraints, evaluated separately, as the geometries
    of D3IL-avoiding are. The first, \emph{tilt}, is a halfspace boundary that slopes by $14^\circ$
    across the corridor in the horizontal plane and is leaned over by $60^\circ$ from the vertical about
    the launch altitude, so that one plane pushes the vehicle sideways and downward at once: at the
    launch altitude a sideways move alone does not clear it on any lane. The second, \emph{hump}, is a
    roof in the vertical plane that rises from the floor at $x=-1.5$\,m to $1.10$\,m at $x=0$ and falls
    back to the floor at $x=1.5$\,m, which the vehicle must climb over; the ceiling of the workspace box
    is raised for it. The models are trained on the demonstrations of a narrower, level corridor and
    never see either boundary, so both are test-time constraints in the same sense as the halfspaces of
    D3IL-avoiding.
  \item[UAV-s-curve] --- the tracking-controller caveat. Two offset corridor segments, joined by a
    diagonal passage, require two sharp changes of direction, and the scene is reported for what those
    turns do to the controller (\autoref{sec:res:uav:controller}), not as a third benchmark.
\end{description}

\begin{figure}[!htbp]
  \centering
  \begin{minipage}{0.78\linewidth}
  \centering
  \begin{minipage}{0.30\linewidth}\centering
    \includegraphics[width=\linewidth]{fig_render_uav_pillars}\\[-0.1em]
    \includegraphics[width=\linewidth]{fig_scene_uav_pillars}\\[0.2em]
    {\footnotesize UAV-pillars}
  \end{minipage}\hfill
  \begin{minipage}{0.30\linewidth}\centering
    \includegraphics[width=\linewidth]{fig_render_uav_corridor}\\[-0.1em]
    \includegraphics[width=\linewidth]{fig_scene_uav_corridor}\\[0.2em]
    {\footnotesize UAV-corridor}
  \end{minipage}\hfill
  \begin{minipage}{0.30\linewidth}\centering
    \includegraphics[width=\linewidth]{fig_render_uav_scurve}\\[-0.1em]
    \includegraphics[width=\linewidth]{fig_scene_uav_scurve}\\[0.2em]
    {\footnotesize UAV-s-curve}
  \end{minipage}
  \end{minipage}
  \caption[The quadrotor scenes]{The three quadrotor scenes. Upper row: the MuJoCo scene of each, seen
  from behind the start, with the Skydio X2 at its start and a reference path in green: on UAV-pillars a
  mapped demonstration of D3IL-avoiding, on the other two the generator's reference route
  (\autoref{sec:setup:data:uav}). Lower row: the same scene files, orthographic, at true proportions; the
  green line at the far end is the finish line, at the end of the walls on UAV-corridor and UAV-s-curve.
  The test-time constraints of UAV-corridor
  are drawn in \autoref{fig:constraints-uav}.}
  \label{fig:env-uav}
\end{figure}

The three scenes use the same two constraint forms as D3IL-avoiding --- halfspaces and circular
keep-out regions --- so the projection solves the same kinds of rows on UAV-pillars as on D3IL-avoiding
(\autoref{fig:constraints-uav}); on UAV-pillars they are the avoiding rows themselves. What differs on the
two scenes with quadrotor demonstrations is that a wall acts only over the span of $x$ on which it exists,
that the corridor's two constraints leave the horizontal plane --- the tilt is one leaned plane, the hump a
roof --- and that every wall, plane, roof and corner disk is
inflated by the rotor reach of $0.31$\,m before it reaches the solver, because the vehicle is a body and
not a point. The tightened variants move each of these boundaries a further $0.025$\,m inward, the margin
of \ac{DPCC}.

\begin{table}[!htbp]
  \caption[Quadrotor scenes]{The constraint set of each quadrotor scene: halfspaces and circular
  keep-out regions (\autoref{sec:method:env:uav}). On UAV-corridor and UAV-s-curve each boundary is
  inflated by the vehicle's rotor reach of $0.31$\,m before it reaches the solver, and the tightened
  variants move it a further $0.025$\,m inward. On UAV-pillars the boundaries are those of D3IL-avoiding
  mapped into the arena, and the vehicle's extent enters through the scale, which sends the rod's radius
  onto the rotor reach.}
  \label{tab:uav-scenes}
  \centering
  \small
  \setlength{\tabcolsep}{8pt}
  \renewcommand{\arraystretch}{1.25}
  \begin{tabularx}{\linewidth}{@{}l>{\raggedright\arraybackslash}p{3.5cm} X@{}}
    \toprule
    Scene & Constraint & Values \\
    \midrule
    UAV-pillars
      & three avoiding geometries, mapped & the halfspaces and the keep-out disk of \emph{top-left-hard}, \emph{top-right-hard} and \emph{both-hard} (\autoref{tab:avoiding-geometries}) at scale $36$: keep-out radius $2.88$\,m; physical pillars of radius $0.90$ and $1.08$\,m in an arena of $21.6\times25.2$\,m \\
    \addlinespace[2pt]
    UAV-corridor
      & two wall halfspaces & inner faces $y=\pm0.95$\,m, active over $x\in[-2,2]$; clear width $1.90$\,m \\
      & \emph{tilt}: the leaned slide & the sloping line from $(-2.0,\,0.95)$ to $(2.0,\,-0.05)$, $14^\circ$ in the horizontal plane, leaned by $-60^\circ$ from the vertical about $z=1.11$\,m: the feasible distance grows by $\tan(-60^\circ)\,(z-1.11)$, so descending gains room; ceiling $1.80$\,m \\
      & \emph{hump}: two roof halfspaces ($x$--$z$) & from $(-1.5,\,0)$ to $(0,\,1.10)$ and from $(0,\,1.10)$ to $(1.5,\,0)$, the plan above both, each active over its own span of $x$; ceiling raised to $2.80$\,m \\
      & wall ends & four disks of radius $0.05$\,m at $x=\pm2$, $y=\pm1$ \\
    \addlinespace[2pt]
    UAV-s-curve
      & four switched wall halfspaces & feasible $y\in[-1.25,-0.35]$ over $x\in[-3,-0.5]$, and $y\in[0.35,1.25]$ over $x\in[0.5,3]$ \\
      & inner corners & two disks of radius $0.05$\,m at $(\mp0.5,\,\mp0.3)$ \\
    \bottomrule
  \end{tabularx}
\end{table}

\autoref{fig:constraints-uav} draws the constraint sets, one column per scene and one panel per set. On
UAV-pillars the three geometries of D3IL-avoiding are mapped into the arena, halfspaces and keep-out disk
at scale $36$, and the six pillars are drawn as the physical obstacles they are; the mapped scale already
carries the vehicle's extent, so no rotor-reach band is added there and the tightening is the mapped
$0.90$\,m. On UAV-corridor the tilt
is drawn twice --- from above ($x$--$y$) at the launch altitude, where the leaned plane cuts the corridor,
and from the side ($x$--$z$) along the corridor, where the lean itself is --- and the hump from the side,
where the roof is; the walls hold under both. UAV-s-curve has one set. Every panel names the plane it is
drawn in. A wall constrains only the span of $x$ over which it exists, which is why the halfspace it
excludes stops short of the panel edge.

\begin{figure}[!htbp]
  \centering
  \includegraphics[width=\linewidth,height=0.70\textheight,keepaspectratio]{fig_constraints_uav}
  \caption[The constraint sets of the quadrotor scenes]{The constraint sets of the quadrotor scenes,
  one column per scene and one panel per set. Solid line: the constraint boundary; grey band: the
  $0.31$\,m rotor reach added to it; dashed: the tightened boundary; faint wash: the rest of the excluded
  region; white: where the plan may go. Left: UAV-pillars, the three geometries of
  \autoref{tab:avoiding-geometries} at scale $36$, the keep-out disk translucent over the pillar it covers,
  the six pillars solid. Middle: UAV-corridor, the tilt from above at the launch altitude and from the
  side, and the hump from the side. Right: UAV-s-curve from above.}
  \label{fig:constraints-uav}
\end{figure}

\FloatBarrier

\section{Datasets}\label{sec:setup:data}

\subsection{Obstacle Avoidance}\label{sec:setup:data:avoiding}

The demonstrations are those of the D3IL obstacle-avoidance task \parencite{jia2024towards}, as used by \ac{DPCC}. Each step
records the commanded and the measured planar end-effector position and the planar increment between
consecutive commanded positions (\autoref{sec:method:env:avoiding}). The dataset holds 96 demonstrations, drawn in \autoref{fig:constraints-avoiding}. None of them satisfies the constraints of top-left-hard, one those of top-right-hard and two those of both-hard (\autoref{tab:avoiding-geometries}), so satisfying the constraints is left to the projection.

\begin{figure}[htpb]
  \centering
  \includegraphics[width=0.90\linewidth]{fig_constraints_avoiding}
  \caption[The demonstrations against the three constraint geometries]{The 96 demonstrations against the
  three constraint geometries of \autoref{tab:avoiding-geometries}, the companion of
  \autoref{fig:expert-aligning} and \autoref{fig:expert-uav}. Shaded: the region excluded by the halfspaces
  and the keep-out disk; dashed: the tightened boundary. A demonstration is drawn red where it crosses
  the geometry of its panel and green where it satisfies it at every recorded step. The number of green
  lines in each panel is the last column of \autoref{tab:avoiding-geometries}.}
  \label{fig:constraints-avoiding}
\end{figure}

\subsection{Vision-Conditioned Alignment}\label{sec:setup:data:aligning}

The demonstrations are those of the D3IL alignment task \parencite{jia2024towards}: its $900$ training
demonstrations, with the two camera streams rendered by MuJoCo. For this thesis each step is laid out as the nine-dimensional plan of
\autoref{sec:method:env:aligning} --- the three-dimensional increment and the commanded and measured
end-effector positions --- and paired with the two images. Cut into plans of eight steps, they give
$168\,274$ training windows.

The demonstrations stand in the same relation to the constraints as the obstacle-avoidance ones.
\autoref{fig:expert-aligning} draws every recorded context --- the $60$ used for training and the $60$
held out --- as the straight push from the box to its target, which is the shortest route a
demonstration could take. The figure is the top view of the table plane, $x$--$y$, the plane the push
happens in; it draws the contexts and not the demonstrated end-effector paths themselves. That segment
crosses the keep-out region in $106$ of the $120$ contexts, and in $119$ of them once the region is
tightened. It crosses the halfspace in none, tightened or not: in every recorded context the box and its
target both lie on the allowed side, so the straight push between them does too. Whether a
demonstration's end effector entered the far corner while working its way around the box cannot be read
from the contexts, and the recorded paths were not tested; the plans learned from them do enter it ---
without projection the end effector crosses the keep-out region in seven of the ten evaluation contexts
and the halfspace in one (\autoref{sec:res:aligning:projection}). The demonstrations were recorded
without either constraint existing, so the task they teach is precisely the one the constraints forbid;
going around them is the projection's work, not something the policy could have imitated.

\begin{figure}[htpb]
  \centering
  \includegraphics[width=0.62\linewidth]{fig_expert_aligning}
  \caption[The recorded contexts against the alignment constraints]{The $120$ recorded contexts of the
  alignment task over the constraint set of \autoref{fig:aligning-contexts}. Each line joins the
  initial box centre (square) to its target centre (circle); red where that straight push crosses the
  circular keep-out region, green where it clears it; no push crosses the halfspace, whose excluded corner
  is shaded. The draw ranges of
  \autoref{sec:setup:tasks:aligning} are visible as the two bands the squares and circles fall in.}
  \label{fig:expert-aligning}
\end{figure}

In all three environments, a random tenth of the training windows is held out and used only to select
the checkpoint (\autoref{tab:train}).

\subsection{Quadrotor Benchmark}\label{sec:setup:data:uav}

Two of the three scenes have quadrotor demonstrations, and one has none. On UAV-pillars the planner is
the one trained on the demonstrations of D3IL-avoiding (\autoref{sec:setup:data:avoiding}) --- human
teleoperation of the manipulator --- and it is not retrained: the scene is that
obstacle field scaled into an arena, and the quadrotor executes the plans the planner produces there
(\autoref{sec:setup:tasks:uav}). The scene therefore compares the execution of the same plans by two
plants, the manipulator of D3IL-avoiding and the quadrotor with its controller.

The demonstrations of UAV-corridor and UAV-s-curve were generated for this thesis as described in
\autoref{sec:method:env:uav}. That section gives the waypoint-and-arc construction, cosine speed profile,
tracking equations, rejection rule and recorded observation--action pair; \autoref{tab:uav-demos} records
only their numerical and scene-specific parameters. The corridor with the sloping halfspace and the roof
reuses the demonstrations and trained models of the level corridor, because both boundaries are test-time
constraints.

\begin{table}[!htbp]
  \caption[Quadrotor demonstrations]{Generation of the quadrotor demonstrations of UAV-corridor and
  UAV-s-curve; UAV-pillars has none and uses the avoiding planner. Shared by the two scenes:
  altitude drawn uniformly from $[0.90,1.30]$\,m per episode, physics step $0.01$\,s, recording every
  third physics step ($0.03$\,s), the nominal gains of the tracking controller, rejection of an episode
  that falls below $0.50$\,m or touches an obstacle on more than the stated fraction of its steps.
  Durations and offsets are drawn uniformly from the stated ranges.}
  \label{tab:uav-demos}
  \centering
  \footnotesize
  \setlength{\tabcolsep}{3pt}
  \renewcommand{\arraystretch}{0.92}
  \begin{tabularx}{\linewidth}{@{}l L L L@{}}
    \toprule
    & UAV-pillars & UAV-corridor & UAV-s-curve \\
    \midrule
    Passage classes & --- (the D3IL-avoiding demonstrations) & three lanes at $y\in\{-0.12,\,0,\,0.12\}$\,m & one \\
    Reference path & --- & straight, $x$ from $-2.8$ to $2.8$\,m & waypoints joined by arcs of radius $\le0.45$\,m \\
    Duration & --- & $[6,10]$\,s & $[16,22]$\,s \\
    Lateral offset & --- & start, $[-0.03,0.03]$\,m, centre lane & whole path, $[-0.04,0.04]$\,m \\
    Maximum collision fraction & --- & $2\,\%$ & $8\,\%$ \\
    Episodes generated & --- & 500 & 500 \\
    Episodes accepted & --- & 500 & 500 \\
    \bottomrule
  \end{tabularx}
\end{table}

These demonstrations are not human. There is no public teleoperated quadrotor dataset for these scenes,
so the data this thesis trains on was \emph{generated} for it: a reference path is laid out analytically
per passage class (\autoref{tab:uav-demos}), and every episode is then \emph{flown in MuJoCo} by the same
cascaded geometric controller the evaluation uses (\autoref{sec:method:deployment}), under the same
physics and at the same time step. What is recorded is the measured state of that flight, not the
reference. An episode is kept only if the flight it produced stays airborne and within its collision
budget (\autoref{tab:uav-demos}). Each
demonstration is therefore a trajectory the plant has been shown to track, and \emph{expert} means that
in this thesis and nothing more: a competent generated flight, not a human one.

Where these demonstrations differ from the two manipulation datasets is that they were generated against
the obstacles they fly among, and \autoref{fig:expert-uav} shows what follows from it. The figure draws
the demonstrations of the two scenes that have their own over their constraint sets, the tightened boundary dashed, and the grey silhouette
on a route is the vehicle to scale at its tightest moment (\autoref{fig:platforms}), for size only. On
UAV-s-curve the
single route, the generator's reference route, satisfies the constraint set the projection later enforces,
tightening included: the generator threads a passage the scene already contains, so the demonstrated
route itself needs no correction. What the projection meets on that scene is the learned plan, which cuts
the second corner of the crossover (\autoref{sec:res:uav:scurve}). The scene is read for what the tracking
controller and the selected per-step configuration make of that plan; it yields no ranking of the
projection methods, because at the selected budget of one evaluation endpoint projection has no guiding
step (\autoref{sec:res:uav:scurve:projection}).

UAV-corridor is built the other way, and it is the way the obstacle-avoidance benchmark is built. The three
lanes clear both walls, but every one of them crosses the tilt --- at the launch altitude, where the leaned
plane cuts the corridor along the same line as before it was leaned, by $0.24$, $0.36$ and $0.47$\,m for the
left, centre and right lane --- and every one of them passes under the hump, whose roof at its peak stands
$1.48$\,m above the floor once inflated while the demonstrations fly between $0.90$ and $1.30$\,m. Both
boundaries are introduced at test time and did not exist when the demonstrations were flown. The scene in
which the constraint is new to the data is exactly the scene in which the projection has something to do.

UAV-pillars is the second kind by construction, and more sharply than the corridor: its demonstrations are
the $96$ of D3IL-avoiding (\autoref{fig:constraints-avoiding}), mapped into the arena at scale $36$ by the
same similarity map as its constraint set, so they are not drawn again in \autoref{fig:expert-uav}, and
its constraint sets are the three test-time geometries of
that task, which $0$, $1$ and $2$ of the $96$ demonstrations satisfy (\autoref{sec:setup:tasks:avoiding});
so on that scene the projector has the same repair to do as on D3IL-avoiding, against the same constraint
sets, which leave five of the six pillars out (\autoref{sec:setup:tasks:uav}).
\autoref{sec:res:uav:pillars} reads it on those terms, on the configurations that win on D3IL-avoiding.

On the two quadrotor-demonstrated scenes that figure draws the references the generator laid out rather
than the flights that tracked them, which is the conservative choice: a constraint crossed by the reference is crossed by construction rather than
by tracking error.

One thing has to be read correctly in that figure and in every UAV-corridor and UAV-s-curve result that
follows. A violation is a property of the \emph{path}, not of the aircraft: what is tested is the line of
vehicle centres against the constraint surfaces, and on those two scenes the surfaces already carry the
$0.31$\,m rotor reach (\autoref{tab:eval}), so the vehicle's own extent is inside the boundary rather than
measured against it. On UAV-pillars no band is added: the scale of $36$ that maps the table into the arena
sends the rod's radius onto the vehicle's radial reach of $0.36$\,m, so the mapped constraint sets carry
the extent already (\autoref{sec:setup:tasks:uav}).
A route turns red when its centre line crosses an inflated surface, and nothing is ever scored by sweeping
the rotor disks through an obstacle. The scenes are built to be flyable on those terms --- the free
channel that remains after inflation is a real one --- so a red route means a plan that had somewhere else
to go, not a vehicle too large for the gap.

\begin{figure}[htpb]
  \centering
  \includegraphics[width=\linewidth,height=0.70\textheight,keepaspectratio]{fig_expert_uav}
  \caption[The demonstrations against the quadrotor constraints]{The demonstrations of the two
  quadrotor-demonstrated scenes over their constraint sets, drawn as in \autoref{fig:constraints-uav}:
  UAV-corridor, the three reference lanes against the tilt from above and the flown altitude band against
  the hump from the side; UAV-s-curve, its single route. UAV-pillars flies the demonstrations of
  D3IL-avoiding (\autoref{fig:constraints-avoiding}) and is not drawn. Green satisfies the constraints,
  red crosses them; the tightened boundary is dashed and the grey silhouette is the vehicle to scale.}
  \label{fig:expert-uav}
\end{figure}

\section{Metrics}\label{sec:setup:metrics}

\subsection{Obstacle Avoidance}\label{sec:setup:metrics:avoiding}

The metrics are those of \ac{DPCC} \parencite[Sec.~6.1]{romer2025diffusion}: the number of control
steps to reach the goal (steps per episode), the success rate of reaching the goal (fraction of episodes),
the success rate of reaching the goal while satisfying the constraints at every step (fraction of
episodes), and the number of control steps on which a constraint is violated (steps per episode).

The first of those is a headline number in both works this chapter measures against, and not a
quantity introduced here: \ac{DPCC} reports the control steps an episode takes beside its success
rates \parencite[Sec.~6.1]{romer2025diffusion}, and the endpoint-projection work, which follows the setup of
\ac{DPCC} on this same D3IL task and counts projection after every sampling step, the projection of \ac{DPCC},
among its baselines, reports the average number of steps the manipulator takes to reach the target on its
constraint-satisfying episodes as one of its three columns, beside the safety rate and the computation time
\parencite[Sec.~VII-A, Table~II]{li2025hardflow}. This thesis compares the two projection methods in the same
way: endpoint projection against the per-step projection of \ac{DPCC} as the baseline, with the control steps
among the measures. A plan that arrives in
fewer control steps is a plan that completes the task sooner, which is why it is reported next to
success rather than folded into it.

To these this thesis adds the time to compute one action (milliseconds, or seconds where a
configuration is slow enough to warrant it), the quantity \ac{DPCC} states for its own
method \parencite[Sec.~6]{romer2025diffusion}. It is measured around the planner at every control step
and covers everything the planner does to produce the action that is executed: generating the $B$
candidate plans, projecting them onto the constraint set, and selecting one. The same measurement is
used in all three environments, and
\autoref{sec:setup:protocol:compute} gives the machine every configuration was measured on.

\subsection{Vision-Conditioned Alignment}\label{sec:setup:metrics:aligning}

Each episode is scored by
\begin{itemize}
  \item the final distance between the box and its target in the table plane (metres), summarised over the
    contexts by its median and its minimum, and, in per cent, the reduction of the median final distance
    relative to the mean initial distance of the ten contexts,
    $100\,(1-\operatorname{median}(d\sidx{final})/\overline{d}\sidx{init})$ with
    $\overline{d}\sidx{init}=0.4530$\,m: $100\,\%$ is a median box brought onto its target and $0\,\%$ a
    median final distance equal to the mean initial one --- not a box left unmoved, since the initial
    distances differ between contexts (their median is $0.4557$\,m);
  \item whether the box was moved at all (a count of contexts);
  \item whether the box ends in position: within $0.018$\,m of its target in the table plane, the position
    tolerance of the task's own success test (a count of contexts); the orientation half of that test is not
    used, for the reason given below;
  \item the number of control steps on which a constraint is violated (steps per episode);
  \item the time to compute one action (milliseconds), as in \autoref{sec:setup:metrics:avoiding}.
\end{itemize}
When generative models are compared, the distance is measured on plans executed without projection, so
that it reflects the generative model alone; when projection methods are compared, it is measured with
projection.

\subsection{Quadrotor Benchmark}\label{sec:setup:metrics:uav}

Each flight is scored by
\begin{itemize}
  \item success: the vehicle crosses the finish line of its scene. This is the success
    criterion of D3IL-avoiding transferred to the quadrotor --- there an episode succeeds when the end effector
    crosses the goal line (\autoref{sec:setup:metrics:avoiding}), here when the vehicle crosses the line
    that closes its route. On UAV-corridor and UAV-s-curve a crossing counts only in controlled flight: the
    flight was not aborted (below), its lowest altitude stayed above $0.2$\,m, and the vehicle collided with
    the scene's walls on no more than $2\,\%$ of its physics steps on UAV-corridor and $8\,\%$ on
    UAV-s-curve, the tolerances with which the demonstrations of each scene were accepted. On UAV-corridor
    and UAV-s-curve the finish line is the end of the walls, $x=2.0$\,m and $x=3.0$\,m, where the walls and
    every test-time constraint end; the crossing is latched the moment the vehicle passes it, the flight flies
    on as before, and the controlled-flight checks above are taken over the whole flight. On
    UAV-corridor the line through the route's goal point, $0.8$\,m further on, is not used: the episode limit was set from the
    demonstrations of the level corridor, and a projected flight that has descended under the tilt or climbed
    over the hump can reach that limit after the constraints have ended while still flying towards the goal.
    The counts at the goal point are printed beside the others where they differ. Every quadrotor
    table and figure of \autoref{ch:results} calls a flight that does so a \emph{success} and prints the
    count of them;
  \item whether it is violation-free: no control step on which the flown path violates a constraint. The
    test is on the path, not on the airframe: on UAV-corridor and UAV-s-curve the surfaces it is tested
    against already carry the $0.31$\,m rotor reach (\autoref{tab:eval}), and on UAV-pillars the mapped
    constraint sets carry the vehicle's extent through the scale (\autoref{sec:setup:tasks:uav});
  \item success with constraint satisfaction: success on a violation-free flight, as on
    D3IL-avoiding;
  \item on UAV-corridor and UAV-s-curve, where almost no flight is free of violations, the constraint is read
    as the number of \emph{violating control steps} per flight, which the tables of those scenes print in place
    of success with constraint satisfaction. Fewer is better, and the cost frontier of UAV-corridor is drawn on
    violating steps against the time to compute one action, among the configurations that reach the finish
    line on every flight; within groups of configurations that violate about equally it is drawn as on
    D3IL-avoiding, on control steps to the finish line against time;
  \item on UAV-pillars, the metrics of D3IL-avoiding (\autoref{sec:setup:metrics:avoiding}) scored on the
    flown path --- success at the mapped goal line, success with constraint satisfaction, violating
    control steps, control steps --- so that a plan's outcome on D3IL-avoiding and on UAV-pillars is read
    in the same units. A flight in which the vehicle collides with a pillar ends there and is not a
    success; the collision is not a constraint violation, since the pillars it concerns are in no
    constraint set, and the tables count it through the success column;
  \item whether the flight was \emph{aborted}: the evaluation ends a flight the moment the vehicle
    inverts --- its body vertical axis points below the horizon --- because a vehicle that has lost
    attitude control does not recover and would otherwise be flown on to the episode limit. The guard
    tests the attitude only, never a collision. An aborted flight is not a success, its outcome columns
    are read as they stood when the guard fired, and the control step at which it fired is recorded.
    The quadrotor tables of \autoref{ch:results} print the count of aborted flights per cell;
  \item the number of control steps (steps per flight), and the final distance to the goal point of the
    route (metres), which describes where a flight that does not succeed ends;
  \item the time to compute one action (milliseconds, \autoref{sec:setup:metrics:avoiding}), split into
    generation and projection.
\end{itemize}

\subsection{Aggregation and Variability}\label{sec:setup:metrics:spread}

Every metric above is measured once per episode and reported as an average over many of them, so a bare
average would hide how much it moves. \autoref{ch:results} therefore prints \emph{mean $\pm$ spread}
wherever a spread exists, and what the spread is over depends on what the result is a claim about.

\begin{itemize}
  \item \textbf{Multi-seed results} --- the D3IL-avoiding tables at the protocol of \ac{DPCC}, and
    UAV-pillars, which follows it --- are reported as the mean over the
    five training seeds and the \emph{sample standard deviation across those seeds}. A claim there is a
    claim about a model and not about one training run, so the variation that matters is how much the
    result moves when the model is trained again. It is also the larger of the two spreads available and
    therefore the conservative one. Each seed is first averaged over the constraint geometries it covers,
    then the mean and the standard deviation are taken over the seeds
    (\autoref{sec:setup:protocol:avoiding}).
  \item \textbf{Single-seed results} --- the one-seed D3IL-avoiding tables, D3IL-aligning, UAV-corridor and
    UAV-s-curve --- carry no seed dimension, so the spread is the sample standard deviation
    \emph{across the episodes, contexts or flights} of the cell, which is stated with each table. Where a
    cell averages several constraint geometries, the within-geometry standard deviations are pooled.
  \item \textbf{Counts are printed as counts.} A result of the form \enquote{$10/10$ flights} or
    \enquote{$20/20$ episodes} is not an average and gets no $\pm$; neither does a median, which is
    reported with the number of contexts it is taken over.
\end{itemize}

\section{Baselines}\label{sec:setup:baselines}

Three generative models are compared with the diffusion model of \ac{DPCC} (\emph{Diffusion} in tables):
\emph{instantaneous-velocity matching} (\emph{FM}) \parencite{lipman2023flow}, which learns the
instantaneous velocity of the probability path; \emph{analytic average-velocity matching}
(\emph{MeanFM}) \parencite{geng2025meanflow}, which learns the average velocity over an interval through
an analytic Jacobian--vector-product target; and \emph{consistency-interpolated average-velocity
matching} (\emph{CI-MeanFM}) \parencite{zhang2025alphaflow}, whose consistency step ratio $\alpha$
interpolates the target between the instantaneous velocity at $\alpha=1$ and a stop-gradient prediction
towards the analytic average-velocity objective as $\alpha\to0$. In this thesis $\alpha$ is annealed from
$1$ and held at a final value
$\alpha_{\mathrm{end}}$, which each result states. Two projection methods are compared: the per-step projection of
\ac{DPCC} and endpoint projection \parencite{li2025hardflow}. All generative models share the temporal
U-Net backbone of the baseline, with $4.0$\,M parameters, unless a table states otherwise. Throughout
all figures and tables, the four models are labelled \emph{Diffusion}, \emph{FM}, \emph{MeanFM}, and
\emph{CI-MeanFM}, respectively; running prose retains the full mechanism names.

\subsection{Obstacle Avoidance}\label{sec:setup:baselines:avoiding}

The baseline is \ac{DPCC} in the configuration its authors report: the diffusion model at a step budget
of $\nfe=20$, per-step projection on tightened constraints, and selection by
cumulative projection cost over $B=4$ candidate plans \parencite[Sec.~6.1]{romer2025diffusion}.
\autoref{tab:target} places the published result next to its reproduction in this thesis. The baseline is
compared on the tightened constraint set, and so is every other configuration of this thesis
(\autoref{sec:setup:tasks:avoiding}): tightening is the margin \ac{DPCC} itself introduces so that a plan
projected onto the shrunk set stays feasible under the tracking and discretisation error of execution
\parencite[Sec.~5.3]{romer2025diffusion}, it is the setting under which \ac{DPCC} reports its results, and it
is the setting that gives the better result. \ac{DPCC} reports the same configuration with and without
tightening \parencite[Table~1 and Fig.~2]{romer2025diffusion}: without it, the rate of success with
constraint satisfaction is $0.49$ at $76.1\pm21.5$ steps with $6.0\pm11.5$ violating steps per episode,
and $0.41$, $0.51$ and $0.49$ across the three selection criteria; with it, $0.98$ at $69.0\pm12.9$ steps
with $0.0\pm0.3$ violating steps, and $0.96$ to $0.99$ across the criteria. The first two rows of
\autoref{tab:target} are those numbers, the third the reproduction of the tightened one. The
reason is the margin: a plan that clears the tightened boundary clears the real one under the tracking
and discretisation error of execution; a plan projected onto the plain boundary has no such margin.

\begin{table}[htpb]
  \caption[The DPCC baseline, published and reproduced]{The DPCC baseline on the obstacle-avoidance
  task: the published result without and with constraint tightening, and the reproduction of the
  tightened configuration in this thesis, on the metrics of \autoref{sec:setup:metrics:avoiding}. S\&C denotes success with constraint satisfaction; time is
  seconds per control step. The spread on the reproduced row is the standard deviation over episodes,
  pooled over the five training seeds within a geometry and then averaged over the three geometries, so it
  carries both the seed-to-seed and the episode-to-episode variation.}
  \label{tab:target}
  \centering
  \small
  \begin{tabularx}{\linewidth}{@{}L c c c c c@{}}
    \toprule
    Source and protocol & Steps & Success & S\&C & \shortstack{Violating\\steps} & s/step \\
    \midrule
    Published, no tightening \parencite[Table~1]{romer2025diffusion}
      & $76.1\pm21.5$ & $0.96$ & $0.49$ & $6.0\pm11.5$ & --- \\
    Published, tightened \parencite[Table~1]{romer2025diffusion}
      & $69.0\pm12.9$ & $1.00$ & $0.98$ & $0.0\pm0.3$ & --- \\
    Reproduced, tightened, protocol of \ac{DPCC}
      & $70.1\pm9.0$ & $1.000$ & $1.000$ & $0.0\pm0.0$ & $0.553$ \\
    \bottomrule
  \end{tabularx}
\end{table}

The reproduction uses the five training seeds and two episodes per seed and geometry that \ac{DPCC} itself
reports on (\autoref{sec:setup:protocol:avoiding}). A twenty-episode evaluation of two of the flow-based
models at the budgets they are operated at exists and agrees with this protocol
(\autoref{app:avoiding-twenty}); the baseline's own twenty-episode run did not complete and is not used.

\subsection{Vision-Conditioned Alignment}\label{sec:setup:baselines:aligning}

No published result exists for the configuration evaluated here, so there is no external baseline to
reproduce, as there is on D3IL-avoiding. The task itself is D3IL's \parencite{jia2024towards}, which
evaluates imitation-learning policies on it from images: two $96\times96$ RGB views, front and in-hand,
each passed through its own ResNet-18 convolutional trunk with spatial-softmax pooling.

What D3IL scores as \emph{success} on this task is a threshold on the final pose of the box, and both
halves of it must hold at once: the box centre within $0.018$\,m of the target centre \emph{and} the box
orientation within $0.048\pi \approx 8.6^{\circ}$ of the target orientation, the second of which is what
makes the coloured sides line up. An episode that meets both counts as $1$ and every other episode as $0$;
the published figure is the share of episodes counted $1$.

That metric is not used here. The planner of this thesis is a different network with a different
training target and output --- a plan of positions that is projected and tracked, not an action applied
directly (\autoref{tab:va-vs-d3il}) --- so the published success rate does not transfer to it. The
alignment task is scored by the final distance of the box to its target
(\autoref{sec:setup:metrics:aligning}), and the D3IL numbers that follow describe how hard the task is.

From images the strongest policy is BESO, a
score-based policy with a history of observations, at $0.359\pm0.031$; next are DDPM-GPT at
$0.316\pm0.026$ and DDPM-ACT, which adds action chunking, at $0.278\pm0.071$. From the state the same task
is solved far more often: VAE-ACT reaches $0.891\pm0.022$ and BESO $0.861\pm0.016$. On obstacle avoidance
from the state BESO reaches $0.950\pm0.015$ \parencite[Table~3]{jia2024towards}. Seen from images, this is a
task that the policies of D3IL finish in roughly a third of episodes, where obstacle avoidance from the
state (D3IL-avoiding) is close to solved.

Those policies are small, and most of what they carry is perception. The \emph{head} --- everything behind
the image encoder --- is a four-layer transformer at an embedding width of $72$ for BESO and DDPM-GPT. For
DDPM-ACT and VAE-ACT it is an encoder--decoder pair at a width of $64$, two layers in the encoder and four
in the decoder; VAE-ACT carries a third two-layer stack at the same width, the variational branch that
produces its latent \parencite[Table~9]{jia2024towards}.

Counted out at those shapes, with the four-fold feed-forward expansion the implementation uses, a
width-$72$ block is $63.0$\,k parameters and the four-layer stack $0.25$\,M. A width-$64$ encoder block is
$49.9$\,k and a decoder block $62.3$\,k --- the difference is the three cross-attention projections --- so
DDPM-ACT's two stacks come to $0.35$\,M and VAE-ACT's three stacks, with its embeddings and output heads,
to $0.47$\,M. None of the four heads reaches half a million parameters.

What carries the weight of an image-conditioned policy there is the encoder, and it is not close: each of
the two camera streams has its own ResNet-18 trunk of roughly $11$\,M parameters, so the perception is
about $22$\,M against a policy of under half a million --- some fifty times the head it feeds. This thesis
takes that encoder over unchanged (\autoref{sec:method:backbone}) and places the $4.0$\,M temporal U-Net
of the state-based benchmark behind it, so the two set-ups differ in the planner and not in how the
images are seen.

The two set-ups also differ in what the network produces and in what happens to it afterwards
(\autoref{tab:va-vs-d3il}), which is the second reason those numbers are a reference and not a baseline. The comparison for the flow-based models is therefore made inside the thesis:
the diffusion model of \ac{DPCC}, trained on the same demonstrations with the same visual encoder,
backbone, horizon and projection, at $\nfe=20$.

\begin{table}[htpb]
  \caption[Alignment set-up against D3IL]{The alignment task as evaluated in D3IL and in this thesis.
  The simulator, the demonstrations and the visual encoder are shared; the policy and the closed loop are
  not.}
  \label{tab:va-vs-d3il}
  \centering
  \small
  \begin{tabularx}{\linewidth}{@{}p{2.6cm} L L@{}}
    \toprule
    & D3IL \parencite{jia2024towards} & This thesis \\
    \midrule
    Network output & actions, executed directly & a plan of $H=8$ steps of action and state \\
    \addlinespace[4pt]
    Network        & MLP, GPT and action-chunking transformer heads, under $1$\,M parameters, behind two ResNet-18 encoders & the temporal U-Net of the state-based benchmark, $4.0$\,M parameters, conditioned on a $128$-dimensional visual latent from the same two encoders (\autoref{sec:method:backbone}) \\
    \addlinespace[4pt]
    Closed loop    & the policy's action is applied & receding horizon: the plan is projected onto the constraint set, its first action is applied, and a new plan is generated at the next control step \\
    \addlinespace[4pt]
    Constraints    & none & the constraint set of \autoref{tab:aligning-constraints}, enforced by the projection \\
    \addlinespace[4pt]
    Scoring        & success rate and entropy of the behaviours & final distance of the box to its target, violations, time per control step (\autoref{sec:setup:metrics:aligning}) \\
    \addlinespace[4pt]
    Reported result & a success rate: at best $0.359 \pm 0.031$ from images, $0.891 \pm 0.022$ from the state & no success rate is defined; an episode is scored by how close the box ends to its target \\
    \bottomrule
  \end{tabularx}
\end{table}
\FloatBarrier

\subsection{Quadrotor Benchmark}\label{sec:setup:baselines:uav}

The quadrotor benchmark is new, and no learned planner has been evaluated on it before. Its closest
precedent is the quadrotor task distributed with MuJoCo MPC \parencite{howell2022predictive}, which flies
the same Skydio X2 model \parencite{zakka2022menagerie}. That task is described here as it is
distributed, because the benchmark of this thesis starts from its vehicle model and departs from it in
everything else.

\paragraph{The original task.} The quadrotor task distributed with MuJoCo MPC
\parencite{howell2022predictive} is the following. The scene is a racing course of eight square gates, each a frame of four
capsules with a $1.5$\,m opening, placed at different headings and heights around a closed loop. The
vehicle starts on the floor. A goal marker leads it around the course: it sits at one of twelve fixed
positions, between $0.75$ and $2.25$\,m above the floor, placed so that flying from one to the next passes
through a gate, and it advances to the next position whenever the vehicle comes within $0.5$\,m of it,
returning to the first after the last. The gates do not appear in the objective; passing through them is
a consequence of where the goal positions are placed.

The controller is Predictive Sampling, the sampling-based predictive controller of MuJoCo MPC
(\autoref{sec:method:mjpc}), and the simulator itself is its model. At each planning instant it perturbs its current control sequence into 32
candidate sequences of the four rotor thrusts, each parameterised by five spline points, rolls every
candidate forward in MuJoCo over a horizon of $0.5$\,s at a time step of $0.01$\,s, and keeps the
cheapest. The cost of a rollout sums four quadratic terms over the horizon: the distance of the vehicle
from the current goal position, weighted $25$; its linear velocity and its angular velocity, each
weighted $1.25$; and the deviation of each rotor thrust from the hover thrust, weighted $10^{-3}$. The
first control of the kept sequence is applied to the rotors, and planning repeats. Nothing is learned,
no demonstrations are used, and no success rate is defined: the task is an interactive demonstration of
the controller.

\paragraph{This benchmark.} The benchmark of this thesis is built on the principle D3IL-avoiding and
\ac{DPCC} are built on \parencite{jia2024towards,romer2025diffusion}: a policy is trained on
demonstrations of a task, constraints the demonstrations do not satisfy are introduced at test time, and
each run is scored both for finishing the task and for satisfying those constraints at every step. It
keeps the vehicle model and the simulator of the original task and
changes the three things that make it a test of constrained generative planning. The plan comes
from a network trained on demonstrations instead of from sampling against the true dynamics; the
environment is a set of walls, pillars and a test-time sloping halfspace that the plan must not violate, instead of
gates that a goal sequence leads through; and the plan is a sequence of positions tracked by the
cascaded geometric controller (\autoref{sec:method:deployment}), instead of rotor thrusts applied
directly. Each flight is scored for success and for constraint satisfaction (\autoref{sec:setup:metrics:uav}).

\section{Training and Evaluation Protocol}\label{sec:setup:protocol}

Every flow-based model is trained once per training seed, and the same checkpoint is then evaluated under
every step budget, projection method and selection rule; the diffusion baseline's budget is fixed with its
noise schedule at training (\autoref{sec:setup:protocol:eval}), so each of its budgets is a separately
trained checkpoint. Otherwise only the evaluation parameters vary between results. This
section gives the machine everything ran on, then the training parameters, then the evaluation ones.

\subsection{Compute Environment}\label{sec:setup:protocol:compute}

All training and evaluation reported in this thesis ran on a single shared GPU node of the
institute's cluster, scheduled with SLURM. The node carries two 16-core AMD EPYC~7282 processors
(32~physical cores, 64~hardware threads) and eight NVIDIA RTX~A5000 accelerators with 24~GB of
device memory each. \textbf{Every job was allocated exactly one GPU}: no result in this thesis comes
from a distributed or multi-GPU run. The full hardware and software inventory, including package
versions, is given in \autoref{app:repro}, \autoref{tab:compute}.

The measurements of \autoref{sec:setup:metrics} were taken on this machine; what the measured times
contain, and what they leave out, is discussed in \autoref{ch:discussion}.

\subsection{Training}\label{sec:setup:protocol:training}

\autoref{tab:train} lists the training parameters and what each controls. The same table covers all three
environments; where a value differs between models it is given per model.

\begin{table}[p]
  \caption[Training parameters]{Training parameters. Values as recorded by the training runs. The effective batch is the batch size times the number of gradient accumulation
  steps. \enquote{Diffusion (\ac{DPCC})} is the baseline of \autoref{tab:target}: on D3IL-avoiding it is the
  training configuration of \ac{DPCC}, taken over unchanged; on the other two environments the same model
  and backbone are trained on that environment's data, with these settings. The visual
  input is two RGB tensors, each $(N,1,3,96,96)$, encoded to a combined $(N,128)$ feature; $N$ is the
  training batch size. The quadrotor column is UAV-corridor and UAV-s-curve; UAV-pillars flies the
  D3IL-avoiding models unchanged.}
  \label{tab:train}
  \centering
  \footnotesize
  \setlength{\tabcolsep}{4pt}
  \renewcommand{\arraystretch}{1.28}
  \begin{tabularx}{\linewidth}{@{}p{2.0cm} L p{1.3cm} p{1.15cm} p{1.5cm} p{1.65cm} p{1.65cm}@{}}
    \toprule
    & & \multicolumn{3}{c}{D3IL-avoiding} & D3IL-aligning & UAV-corridor, UAV-s-curve \\
    \cmidrule(lr){3-5}
    Parameter & Meaning & Diffusion (\ac{DPCC}) & FM & MeanFM, CI-MeanFM & all models & all models \\
    \midrule
    Plan horizon $H$ & steps in one generated plan & 8 & 8 & 8 & 8 & 8 \\
    Plan width $d$ & action and state channels per step & 6 & 6 & 6 & 9 & 12 \\
    Plan input/output & model tensor $(N,H,d)$ & $(N,8,6)$ & $(N,8,6)$ & $(N,8,6)$ & $(N,8,9)$ & $(N,8,12)$ \\
    State input & current-state tensor & $(N,4)$ & $(N,4)$ & $(N,4)$ & $(N,6)$ & $(N,9)$ \\
    Visual input & image streams; exact shape in caption & --- & --- & --- & two RGB views & --- \\
    Training steps & optimiser updates & $10^{5}$ & $10^{5}$ & $10^{5}$ & $10^{5}$ & $10^{5}$ \\
    Batch size $\times$ accumulation & plans per update & $8\times2$ & $8\times2$ & $32\times2$ & $64\times2$ & $8\times2$ \\
    Learning rate & Adam step size & $10^{-4}$ & $10^{-4}$ & $5\cdot10^{-4}$ & $2\cdot10^{-4}$ & $10^{-4}$ \\
    Weights evaluated & raw weights, or their exponential moving average (decay $0.995$) & raw & raw & moving average & raw & raw \\
    Training-time sampling & \textsc{u}: discrete uniform; \textsc{b}: reflected Beta; \textsc{ln}: logit-normal pair. See \autoref{app:training-time-laws} & \textsc{u} & \textsc{b} & \textsc{ln} & by family: \textsc{u}, \textsc{b}, \textsc{ln} & by family: \textsc{u}, \textsc{b}, \textsc{ln} \\
    Checkpoint & which saved weights are evaluated & \multicolumn{5}{p{7.3cm}}{lowest loss on the held-out $10\,\%$ of windows} \\
    \bottomrule
  \end{tabularx}
\end{table}
The generated plan has tensor shape $(N,H,d)$, with training batch size $N$ and $H=8$; its model output has
the same shape. In D3IL-avoiding, the state anchor is $(N,4)$ and the plan is $(N,8,6)$: a planar
position increment and the commanded and measured planar positions. In D3IL-aligning, the state anchor
is $(N,6)$ and the plan is $(N,8,9)$: a spatial position increment and the commanded and measured
spatial positions. Each of its two RGB cameras supplies one $(N,1,3,96,96)$ image tensor; the
concatenated visual feature has shape $(N,128)$. On UAV-corridor and UAV-s-curve, the state anchor is $(N,9)$ and
the plan is $(N,8,12)$: a spatial position increment, commanded position, measured position and
measured velocity; UAV-pillars flies the D3IL-avoiding models and their $(N,8,6)$ plans unchanged
(\autoref{sec:setup:tasks:uav}). The embodiment layouts are introduced in \autoref{sec:method:deployment}.
For the \textsc{b} entry the actual transport time is $1-\operatorname{Beta}(1.5,1)$; for
\textsc{ln} the recorded parameters are $P\sidx{mean}=-0.4$ and $P\sidx{std}=1.0$. Their exact
distributions are given in \autoref{app:training-time-laws}.

The optimiser settings in that table are not uniform across the models, and the reason is the training
target. The four objectives fit different quantities --- a noise residual, an instantaneous velocity, and
an average velocity over an interval --- and a gradient of one is not a gradient of another: the scale of
the loss, and how much it varies between two draws of the transport time, differ by construction. Each
model was therefore given the batch size, step size and weight averaging under which its own loss curve
converged, rather than a single setting imposed on all four. On D3IL-avoiding this makes the two
average-velocity models the ones trained with four times the batch and five times the learning rate of
instantaneous-velocity matching and the diffusion baseline, and the only ones evaluated with the
exponential moving average of their weights; the logit-normal transport-time sampler in their column is
the same kind of choice, made for the same reason.

What is held fixed is what the comparison rests on. All four models share the backbone, its parameter
count, the dataset, the plan horizon and the projection they are read through, and each is evaluated at
the step budgets stated with its results. What varies is how each objective was brought to convergence,
and that recipe is part of what is compared: the results are those of the trained configurations as
\autoref{tab:train} states them, not of the objectives in isolation.

\subsection{Evaluation}\label{sec:setup:protocol:eval}

\autoref{tab:eval} lists the evaluation parameters, which are the quantities \autoref{ch:results} varies or
holds fixed. The subsections after it state, per environment, how many training seeds and how many episodes
a result rests on. D3IL-avoiding is evaluated over five training seeds, so a result there carries the variation across seeds,
and so is UAV-pillars, which flies its models; D3IL-aligning, UAV-corridor and UAV-s-curve are evaluated
with one training seed, within the storage and compute available to this work.

\begin{table}[htpb]
  \caption[Evaluation parameters]{Evaluation parameters, and the value each takes in each environment.
  A result that departs from the default states so. UAV-pillars flies the planner of D3IL-avoiding and
  takes that column's values, in the table's coordinates; only its low-level control is the quadrotor's.}
  \label{tab:eval}
  \centering
  \small
  \begin{tabularx}{\linewidth}{@{}p{2.9cm} L p{1.9cm} p{1.9cm} p{2.2cm}@{}}
    \toprule
    Parameter & Meaning & D3IL-avoiding & D3IL-aligning & UAV-corridor, UAV-s-curve \\
    \midrule
    Step budget $\nfe$ & sampling steps per plan & per result & per result & per result \\
    Candidate plans $B$ & plans generated in parallel per control step & 4 & 4 & 4 \\
    Selection rule & which of the $B$ plans is executed & per result & per result & per result \\
    Projection method & none, per-step or endpoint & per result & per result & per result \\
    Projected position & coordinates constrained in the plan & $x,y$ (2D) & $x,y,z$ (3D) & $x,y,z$ (3D) \\
    Activation threshold $\eta$ & share of the sampling steps the projection runs on & $0.5$; $1.0$ at the budget floor & lowered with the budget & $0.5$ \\
    Tightening $\delta$ & safety margin, in metres, that every constraint surface is shrunk by; the
      value \ac{DPCC} introduces, and the only quantity the tightened variants change & $0.025$ & $0.03$ & $0.025$ \\
    Body inflation & separate offset, in metres, moving every surface outward so the \emph{vehicle}
      clears it and not only its centre point; the manipulator's end effector is constrained as a point
      and carries none & none & none & $0.31$ \\
    Replanning & how often a new plan is generated & every step & every step & every step \\
    Episode limit & control steps before a run is stopped & 200 & 400 & 396 / 871 \\
    Low-level control & how a setpoint reaches the actuators & \multicolumn{2}{p{3.9cm}}{IK, then joint-space PD} & cascaded geometric tracking control \\
    \bottomrule
  \end{tabularx}
\end{table}

Several of those parameters enter how a result should be read. The projector acts on planar position in
D3IL-avoiding and UAV-pillars, and on spatial position in D3IL-aligning, UAV-corridor and UAV-s-curve;
individual geometric predicates may use only some of those spatial coordinates.

The \emph{step budget} counts the sampling steps of a plan, and it is not the same quantity in the two
families: for the flow-based models it is the number of ODE solver steps, chosen at inference, and for
the diffusion baseline the number of denoising steps, fixed at 20 when the model is trained. Without
endpoint projection it is also the number of network evaluations per plan. Under endpoint projection each
guiding step adds one further evaluation, the lookahead from the predicted state whose endpoint is then
projected, so a plan there costs
$\nfe + n\sidx{guide}$ evaluations (\autoref{sec:setup:protocol:guiding}); the tables name the sampling
steps $\nfe$ throughout.
The \emph{selection rule} is one of the three of \autoref{sec:method:selection} --- random, cumulative
projection cost, or temporal consistency --- and the \emph{projection method} one of the three of
\autoref{sec:method:proj}; both vary per result and neither is ever averaged over. \emph{Low-level
control} means inverse kinematics followed by a joint-space PD controller on the manipulator, and on the
quadrotor the cascaded geometric tracker of \autoref{sec:method:deployment}, which MuJoCo MPC replaces in
\autoref{sec:res:uav:controller} and nowhere else.

An \emph{episode} is one run of the task from its initial state to its end: the environment is reset, the
planner replans at every control step and executes the first action of its plan, and the run ends when the
goal is reached --- the goal line crossed, the box delivered, the finish line of the scene flown through
--- or when the \emph{episode limit} is reached. That limit is a hard cap on how long a single run may
go on, and nothing else. A planner that has lost the task would otherwise replan forever, and every control step of the
quadrotor scenes costs a projection solve, so a run that cannot finish is stopped and scored unsuccessful.
The three limits differ because the tasks are of different lengths: the quadrotor limits are set at $1.2$
times the longest demonstration of each scene, which is why the three scenes have three different values.

The \emph{activation threshold} is not held at one value either. It selects the last $\eta\nfe$ sampling
steps, so the cost of a projected plan is the cost of one solve times that count, and at the budgets
D3IL-aligning is run at the default of $0.5$ would mean solving ten programs for every one of the four
candidate plans at every one of up to $400$ control steps. That is not affordable at the scale of this
evaluation, and it is also not necessary: what endpoint projection needs is not a fixed \emph{share} of
the sampling steps but a fixed \emph{number} of steps left to guide. The threshold is therefore lowered as
the budget rises, so that the two alignment budgets are compared at the same number of guiding steps
rather than at the same share. The one movement in the other direction --- $\eta=1.0$ on D3IL-avoiding ---
is the same argument at the opposite end: at the budget floor half of the sampling steps is one step, and
one step is the terminal one, so the threshold is raised until the method has anything to act on at all.
The two families select the same steps except when $\eta\nfe$ is a whole number below $\nfe$: the
flow-based samplers project on the last $\lceil\eta\nfe\rceil$ steps, while the diffusion baseline's gate
includes the step at the threshold itself and projects on one more --- eleven rather than ten at $\nfe=20$
and $\eta=0.5$; at $\eta=1$ both project on every step.
\autoref{sec:setup:protocol:guiding} works that arithmetic through.

The \emph{tightening} margin is $0.025$\,m on D3IL-avoiding and the quadrotor scenes, the value \ac{DPCC}
introduces, and $0.03$\,m on D3IL-aligning (\autoref{tab:aligning-constraints}). It
shrinks the feasible set, so that a plan which satisfies a constraint numerically still has room left to be
tracked, and it is the one quantity the tightened variants of a constraint set change.

\emph{Body inflation} is a second offset and a different quantity, and only the quadrotor carries it. A
constraint is written for a point, and the manipulator's end effector is treated as one; a quadrotor is
not a point, so every surface of the quadrotor scenes is moved outward before the projection ever sees it
--- by the vehicle's rotor reach of $0.31$\,m. The configuration also defines a planning pad of $0.02$\,m
on top of that reach, but every scene evaluated here switches the pad off for the planner, so the planner
and the violation scorer are offset by the same $0.31$\,m. The $0.31$\,m is the
largest per-axis distance from the centre of mass to the rim of a rotor disk at zero yaw: the four rotors
sit at $(\pm 0.14, \pm 0.18)$\,m and each disk has a radius of $0.13$\,m, so the reach is
$0.18+0.13=0.31$\,m across the vehicle and $0.14+0.13=0.27$\,m along it, and the larger of the two is
taken. The two offsets are kept apart in the configuration and are kept apart here, because they answer
different questions: $0.025$ is margin, $0.31$ is geometry.

\begin{table}[htpb]
  \caption[Evaluation protocol per environment]{What one reported result rests on, per environment.
  D3IL-avoiding follows the protocol of \ac{DPCC}. Results are averaged per geometry first and then
  across geometries, so a geometry covered by fewer seeds does not carry more than its share.}
  \label{tab:protocol}
  \centering
  \small
  \setlength{\tabcolsep}{6pt}
  \renewcommand{\arraystretch}{1.2}
  \begin{tabularx}{\linewidth}{@{}l c L X@{}}
    \toprule
    Environment & Training seeds & Episodes per configuration & Averaged over \\
    \midrule
    D3IL-avoiding & 5
      & $5\times2 = 10$ per geometry
      & the three constraint geometries \\
    D3IL-aligning & 1
      & 10 contexts $\times$ 1 episode
      & the ten contexts, compared context by context \\
    UAV-pillars & 5 & $5\times2 = 10$ per geometry & the three constraint geometries, as on D3IL-avoiding \\
    UAV-corridor & 1 & 10 flights, spread over the routes L, C and R & the ten flights \\
    UAV-s-curve & 1 & 10 flights on its single route & the ten flights \\
    \bottomrule
  \end{tabularx}
\end{table}

\subsubsection{Obstacle Avoidance}\label{sec:setup:protocol:avoiding}

\ac{DPCC} averages over \enquote{five training seeds and ten test seeds for each constraint set
formulation} \parencite[Sec.~6.1]{romer2025diffusion}; its released evaluation runs two episodes for each of
five training seeds. This thesis uses those five seeds and two episodes per seed for its main comparison
(\autoref{tab:protocol}). Every unmarked result in \autoref{sec:res:state} carries all five seeds and all
three geometries. The separate twenty-episode evaluation is reported in
\autoref{app:avoiding-twenty}.

The instantaneous-velocity-matching results were produced by the sampler inherited from the baseline, so
that the initial noise, the candidate fan and the projection call are the baseline's own
(\autoref{sec:method:engine}). That initial noise has standard deviation $0.5$, as in the baseline,
while the model was trained from noise of standard deviation $1$; the two average-velocity models are
sampled from noise of standard deviation $1$, the scale they were trained at. The same holds on
D3IL-aligning and on the quadrotor scenes.

\subsubsection{Vision-Conditioned Alignment}\label{sec:setup:protocol:aligning}

The evaluation resets the environment to a fixed, enumerated set of contexts --- the initial box and target
poses --- so every model and projection method starts from the same ten situations
(\autoref{tab:protocol}). Because the contexts are shared, two configurations are compared context by
context: for each context, which of the two leaves the box closer to its target. A result is the number of
contexts in which one configuration is ahead, together with the median final distance over the ten.

\subsubsection{Quadrotor Benchmark}\label{sec:setup:protocol:uav}

Each scene is reported on its own. UAV-pillars is evaluated at the protocol of D3IL-avoiding --- five
training seeds, two episodes per constraint geometry, averaged per geometry and then over the three
(\autoref{sec:setup:protocol:avoiding}) --- for the configurations \autoref{ch:results} selects there.
The vehicle flies there at a fixed altitude of $1.0$\,m; each commanded position is tracked for one
control step of $1$\,s, $100$ physics steps of $0.01$\,s, with the reference moving towards it at no more
than $1$\,m/s and without velocity feed-forward, and a flight that touches a pillar ends there. UAV-corridor is evaluated with ten flights per configuration, spread
over its left, centre and right routes, on each of its two constraints separately; its projection constrains the commanded position in addition to the measured one and leaves the
action unconstrained, a configuration the other two scenes do not use. UAV-s-curve is evaluated with ten
flights per configuration on its single route: the flow-based models at one, two and twenty network
evaluations and the diffusion baseline at its own twenty, without projection; the configuration that
crosses the finish line most often, under the per-step projector of \ac{DPCC} with each of its three
selection rules, since at the one evaluation that configuration uses endpoint projection has no guiding
step (\autoref{sec:res:uav:scurve:projection}); and its plans under the two tracking controllers. Every flight is tracked by the cascaded geometric
controller of \autoref{sec:method:geometric}; the controller comparison of
\autoref{sec:res:uav:controller} replaces it by MuJoCo MPC and says so.

\subsection{Guiding Steps of Endpoint Projection}\label{sec:setup:protocol:guiding}

Two of the evaluation parameters above interact, and their product decides whether a configuration can
carry a statement about endpoint projection at all. One worked example is enough to read every projected
result in \autoref{ch:results}.

Endpoint projection runs on the sampling steps the activation threshold selects, which are the last
$\lceil \eta\nfe \rceil$ of the $\nfe$ steps. Of those, the terminal step is not a guiding one: there the
predicted endpoint \emph{is} the finished sample, no later step is left for the network to react to, and
what runs is a projection applied to a finished plan --- which is what per-step projection does at that
step too. The steps that remain are the guiding ones:
\begin{equation*}
  n\sidx{guide}(\nfe,\eta) \;=\; \lceil \eta\nfe \rceil - 1 ,
\end{equation*}
as described in \autoref{sec:method:degenerate}.

Read it on the two settings that matter most in \autoref{ch:results}. On D3IL-aligning at
$\nfe=20$ with $\eta=0.2$, the projection runs on $\lceil 4 \rceil = 4$ sampling steps and three of them
are guiding ones, so the reported endpoint-projection rows there are the method doing its own work. On
D3IL-avoiding at the budget the benchmark is actually run at, $\nfe=2$ with the default $\eta=0.5$, the
projection runs on $\lceil 1 \rceil = 1$ sampling step and that step is the terminal one --- so
$n\sidx{guide}=0$ and endpoint projection reduces to a terminal projection of the sampled plan. On
instantaneous-velocity matching a run at $\nfe=2$ that handed both projectors the same candidate plans
produced bit-identical rollouts; on the two average-velocity models the two arms differ even there, because
the endpoint sampler integrates a different field (below), so only on instantaneous-velocity matching is the
difference between the methods nil at that budget. The same arithmetic gives zero guiding steps at $\nfe=1$
for every threshold below $1.0$.

Under endpoint projection the sampler is not the plain one on the two average-velocity models, whether or
not a step is a guiding one: at every transport step it queries the network at a zero interval, where the
plain sampler queries it at the step width $1/\nfe$. On those models a difference between endpoint and
per-step projection therefore belongs to the sampler and the point of projection together, at every budget.

Three rules follow from that, and they hold for every result in \autoref{ch:results}.

\begin{itemize}
  \item A configuration with $n\sidx{guide}=0$ is reported separately and never averaged with
    configurations that have guiding steps.
  \item The count is not assumed. The evaluation records how many programs the projection solved per
    plan, and that number is compared against $\lceil \eta\nfe \rceil$ before a result is reported.
  \item The threshold is stated with every projected result that departs from the default of
    \autoref{tab:eval}.
\end{itemize}

Which configurations the first rule excludes, and how closely the recorded program counts follow the
prediction, is reported with the results (\autoref{sec:res:constraints:degenerate}).
